\subsection{Идеальный упруго-пластический материал Прандтля}

Пружина $\mu$ последовательно с элементом Сен-Венана с пределом $\tau_s$.

\inputfig{figures/fig-prandtl}

\[
  \Sdev = \Sdev_1 = \Sdev_2 ,
  \qquad
  \dev = \dev_1 + \dev_2 ,
  \qquad
  \Sdev_1 = 2\mu\,\dev_1 ,
\]
\[
  \begin{cases}
    N(\Sdev_2) < \tau_s: & \dt{\dev}_2 = 0, \\[0.3em]
    N(\Sdev_2) = \tau_s: & \Sdev_2 = \lambda\,\dt{\dev}_2 .
  \end{cases}
\]
Так как $\dev = \dev_1 + \dev_2$, то $\dt{\dev} = \dt{\dev}_1 + \dt{\dev}_2$, и
полное соотношение распадается на два режима:
\[
  \begin{cases}
    N(\Sdev) < \tau_s: & \dt{\dev} = \dfrac{\dt{\Sdev}}{2\mu}
      \ \Longleftrightarrow\ \dev = \dfrac{\Sdev}{2\mu} , \\[0.9em]
    N(\Sdev) = \tau_s: & \dt{\dev} = \dfrac{\dt{\Sdev}}{2\mu} + \dt{\dev}_2 .
  \end{cases}
\]
Первый режим — чистая упругость, во втором добавляется течение.

\subsubsection*{Пропорциональное деформирование}

Пусть тензорное направление фиксировано, меняется только интенсивность:
\[
  \dev = \dev_0\,\gamma(t),
  \qquad \dev_0 = \mathrm{const},
  \qquad N(\dev_0) = 1 ,
\]
тогда
\[
  N(\dev) = \abs{\gamma}\,N(\dev_0) = \abs{\gamma} ,
\]
$\gamma$ — \textbf{интенсивность сдвиговой деформации}. Напряжение направлено
так же:
\[
  \Sdev = \dev_0\,\tau^*(t),
  \qquad
  N(\Sdev) = \abs{\tau^*(t)} = \tau
  \quad \text{— интенсивность касательных напряжений.}
\]

\subsubsection*{Разбор случаев}

\textbf{1) $\abs{\tau^*} < \tau_s$.} Элемент Сен-Венана заперт,
$\dt\gamma_2 = 0$:
\[
  \dt\gamma = \frac{\dt\tau^*}{2\mu}
  \quad\Longrightarrow\quad
  \gamma = \frac{\tau^*}{2\mu} + C ,
\]
а из начального условия $t = 0$: $\gamma = 0$, $\tau^* = 0$ следует $C = 0$,
то есть
\[
  \tau^* = 2\mu\gamma .
\]

\medskip
\textbf{2) $\tau^* = \tau_s$, $\dt\gamma > 0$ — активное течение.} Напряжение
держится на пределе, растёт пластическая часть. Множитель $\lambda$ находится
из $\tau_s = \lambda\,\dt\gamma$:
\[
  \lambda = \frac{\tau_s}{\dt\gamma} > 0 \quad\text{— допустимо.}
\]
Значит течение возможно: $\gamma$ растёт при $\tau^* = \tau_s = \mathrm{const}$,
и $\gamma_2 = \gamma - \dfrac{\tau_s}{2\mu}$.

\medskip
\textbf{3) $\dt\gamma < 0$ — разгрузка.} Формально
\[
  \lambda = \frac{\tau_s}{\dt\gamma} < 0 ,
\]
что невозможно ($\lambda \ge 0$ по термодинамике). Значит течения нет:
пластическая деформация замораживается на достигнутом уровне
$\gamma_2 = \gamma_{\max} - \dfrac{\tau_s}{2\mu}$, и отклик снова упругий —
прямая с тем же наклоном $2\mu$:
\[
  \tau^* - \tau_s = 2\mu\,(\gamma - \gamma_{\max}) .
\]
Разгрузка идёт по этой прямой до $\tau^* = -\tau_s$, после чего начинается
течение в обратную сторону. При знакопеременном нагружении диаграмма
замыкается в параллелограмм — \textbf{петлю гистерезиса}; её площадь равна
энергии, рассеянной за цикл.

\inputfig{figures/fig-prandtl-cycle}

\subsection{Упруго-пластический материал с упрочнением}

Полное название на доске: «идеальный упруго-пластический материал с линейным
упрочнением».

В модели Прандтля после достижения предела напряжение не растёт вовсе. Чтобы
получить упрочнение, параллельно элементу Сен-Венана ставится вторая пружина
$\mu_1$, а вся связка соединяется последовательно с пружиной $\mu$.

\inputfig{figures/fig-hardening}

Пока трение заперто, блок не деформируется и работает только пружина $\mu$:
\[
  \tau^* = 2\mu\gamma ,
  \qquad \abs{\tau^*} < \tau_s .
\]
После начала течения напряжение в блоке складывается из предела трения и
отклика пружины $\mu_1$:
\[
  \tau^* = \tau_s + 2\mu_1\gamma_2 ,
  \qquad
  \gamma = \frac{\tau^*}{2\mu} + \gamma_2 ,
\]
откуда
\[
  \gamma = \frac{\tau^*}{2\mu} + \frac{\tau^* - \tau_s}{2\mu_1}
  \quad\Longrightarrow\quad
  \boxed{\ \frac{d\tau^*}{d\gamma}
    = \frac{2\mu\mu_1}{\mu + \mu_1}
    = \frac{2}{\dfrac{1}{\mu} + \dfrac{1}{\mu_1}}\ }
\]
— после предела диаграмма идёт прямой с меньшим наклоном: \textbf{линейное
упрочнение}. Предельные случаи: $\mu_1 \to 0$ даёт модель Прандтля (площадка),
$\mu_1 \to \infty$ — чисто упругое тело.

\subsection{Вязко-упруго-пластический материал}

Если вместо второй пружины параллельно элементу трения поставить демпфер
$\eta$, получается вязкопластическая модель: за пределом течения сопротивление
даёт вязкость.

\inputfig{figures/fig-bingham}

\[
  \abs{\tau^*} \le \tau_s: \quad \dt\gamma_2 = 0 ;
  \qquad
  \abs{\tau^*} > \tau_s: \quad
  2\eta\,\dt\gamma_2 = \tau^* - \tau_s\,\sgn\tau^* ,
\]
то есть скорость пластического течения пропорциональна превышению напряжения
над пределом. В тензорном виде
\[
  \dt{\dev}_2
  = \frac{1}{2\eta}\left(1 - \frac{\tau_s}{N(\Sdev)}\right)\Sdev ,
  \qquad N(\Sdev) > \tau_s ,
\]
а полная скорость деформации добавляет упругую часть пружины $\mu$:
\[
  \dt\gamma = \frac{\dt\tau^*}{2\mu} + \dt\gamma_2 .
\]

\begin{warn}
Подпись под схемой в тетради читается как «Бингама» — сверить полное название
модели (Бингама–Шведова?). Сами соотношения восстановлены по виду схемы.
\end{warn}
